\SourceLeaf{133}
schwächen, ja ganz verdunkeln können, ohne darum der Richtigkeit des Satzes den geringsten Eintrag zu thun, dass nach Mittelzahlen aus hinreichend langen Perioden der doppelten Interessentenzahl auch die doppelte Witwenzahl folgen muss. Aber, aus jener Ursache, kann es geschehen, dass bei einer kleinen Gesellschaft die verhältnissmässig vergrösserte Witwenzahl länger ausbleibt, als bei einer grossen; sie kann aber eben so gut auch viel früher eintreten. Es kann, bei einer kleinen Gesellschaft, sich treffen, dass während einer beträchtlichen Zahl von Jahren nach der Vergrösserung der Interessentenzahl die Witwenzahl nur eine ganz unbedeutende Zunahme. zeigt, fast stationär bleibt, ja selbst einmal wieder etwas zurückgeht, was aber im Grunde nichts weniger als wünschenswerth sein würde, falls sich dadurch die Administration in eine trügerische Sicherheit einwiegen liesse, und im Vertrauen auf den augenblicklich noch im Steigen begriffenen Vermögenszustand noch \textit{Erhöhung} der Pension verfügte, zu einer Zeit, wo eine gründliche weiter als auf den nächsten Tag sehende Erwägung vielleicht schon die Nothwendigkeit einer \textit{Beschränkung} erkannt haben würde. Denn das bedarf keines Beweises, dass nothwendig werdende Beschränkungen desto grösser ausfallen müssen, je länger man sie verschoben hatte.

Von dem, was über reine Benefickassen gesagt ist, lässt sich nun leicht die Anwendung auf solche machen, die zwischen jenen und den sich durch die Beiträge ganz selbst erhaltenden stehen. Eine solche gemischte Kasse ist die Professorenwitwenkasse, obwohl sie wegen der Geringfügigkeit der Beiträge jenen viel näher steht als diesen. Auf den Grund jährlicher Beiträge von 10 Thl. würde, wie aus den in der zweiten Abtheilung zu erörternden Rechnungen folgt, den Hinterbliebenen der Interessenten höchstens eine Pension von 44 Thl. oder von 48 Thl. gewährt werden können, je nachdem der Zinsfuss von \(3\tfrac12\) oder von 4 Procent vorausgesetzt wird, und hiebei ist noch nichts wegen möglicher Verluste, und wegen Administrations- und anderer Kosten in Abzug gebracht. Was darüber gewährt wird, also nach dem seit 1835 bestehenden Pensionssatze jährlich 202 bis 206 Thl., ist wie der Ausfluss eines reinen Beneficiums zu betrachten, und es gilt davon, rücksichtlich der Wirkungen der steigenden Interessentenzahl ganz dasselbe, was oben in Betreff solcher Kassen entwickelt ist.

Hiedurch erscheint nun allerdings der Umstand, dass die Anzahl der Theilnehmer an unsrer Witwenkasse jetzt um die Hälfte grösser ist, als sie durchschnittlich vor 20 bis 30 Jahren war, in schwerer Wichtigkeit. Um jedoch diese gehörig würdigen zu können, muss zugleich wohl erwogen werden, dass die in der letzten Zeit so gross gewordene Witwenzahl oder richtiger Pensionenzahl (nach dem \textit{Durchschnitt} der letzten \textit{acht} Jahre \(=20\)) ganz und gar nicht Folge der jetzigen grossen Zahl der Theilnehmer ist, sondern eben so gross sein würde, wenn auch die Zahl der Theilnehmer nicht so sehr vermehrt wäre: es erhellet dies aus dem Umstande, dass die Ehemänner derjenigen Witwen, welche in den letzten acht Jahren den Bestand gebildet haben (resp. Väter der Pension genossen habenden Waisen) fast sämmtlich schon vor dem Steigen der Interessentenzahl, ja meistens schon sehr lange vor diesem Steigen der Gesellschaft angehört haben. Es muss vielmehr die jedesmalige Witwenzahl, in einer Gesellschaft, deren Umfang im Steigen ist, nicht mit der gleichzeitigen Zahl der Theilnehmer, sondern mit derjenigen zusammengestellt werden, welche mehrere Decennien früher Statt gefunden hat. Hiernach liegt nun aber folgende Schlussfolge sehr nahe: Eben so gut, wie aus dem frühern Zustande der Gesellschaft, deren Interessentenzahl vor 20 bis 30 Jahren zwischen 31 und 38 auf und ab schwankte, jetzt eine durchschnittliche Witwenzahl von 20 hervorgegangen ist, wird ganz füglich, wiederum nach einigen Decennien, aus dem jetzigen Umfang der Gesellschaft — von 51 Interessenten — eine Witwenzahl von 30 erwachsen können, und zwar ohne alle Gewähr, dass diese Zahl ein unübersteigliches Maximum sei. Es wird damit nicht gesagt, dass dies gewiss wirklich geschehen \textit{werde}, sondern nur, dass nach den bisherigen Präcedentien es geschehen könne, ohne dass man es gerade wie etwas Ausserordentliches betrachten dürfte; jedenfalls zeigt schon ein solcher roher Ueberschlag, dass die Witwenkasse in den möglichen Wechselfällen ein viel höheres Spiel spielt, als bisher geglaubt sein mag.
